% New EF21-SignMuon obstruction, independent of the fixed-step construction.
% Exact certificate: research/ef21_signmuon_vanishing_steps.py (stdlib only).
\subsection{EF21-SignMuon: decreasing steps and finite budgets}
\label{add:ef_vanishing}

The fixed-step construction scales its gradient floor with the learning
rate. The following example instead keeps the smoothness, initial gap,
dimension, and gradient floor uniform. Failure means nonstationarity;
the objective need not increase without bound.

\begin{lemma}[\textbf{An invisible gradient component}]
\label{add:lem:invisible}
There are $A,B\in\mathbb R^{8\times8}$ and
$S\in\{\pm1\}^{8\times8}$ such that
\begin{gather*}
 \|A\|_F=1,\quad \|B\|_F^2=7,\quad
 \langle A,B\rangle=\langle A,S\rangle=0,\\
 \operatorname{polar}(aA+bB)=A+\operatorname{sign}(b)B
 \quad(a>0,\ b\ne0),\\
 \operatorname{sign}(A+B)=S,\qquad
 \operatorname{sign}(A-B)=-S.
\end{gather*}
Every entry of $A+B$ and $B-A$ has magnitude greater than $1/50$.
Put $h=\langle B,S\rangle/64>0$ and $\kappa=64h^2\le8$.
For alternating targets $D_t=A+(-1)^{t-1}B$, scaled-sign EF21,
initialized at $E_0=0$, produces
\[
 E_t=E_{t-1}+C(D_t-E_{t-1})=(-1)^{t-1}hS.
\]
In particular, $\langle A,E_t\rangle=0$ for every $t$.
\end{lemma}

The algebraic certificate below supplies the matrices. For the recursion,
the two strict sign identities imply
$S_{ij}B_{ij}>|A_{ij}|$ entrywise. The first message is $hS$.
If $E_{t-1}=-shS$ and $D_t=A+sB$, their residual has signs $sS$
and mean modulus $2h$, so $E_t=shS$.

\begin{theorem}[\textbf{Uniform failure of stationarity}]
\label{add:thm:ef_uniform}
For every $b\in\mathbb R$, put
$a_b=\sqrt{1+1/[7(1+b^2)]}$ and define the convex quadratic
\begin{equation}\label{add:eq:quadratic}
 f_b(X)=\frac{(a_b+\langle A,X\rangle)^2}{2}
       +\frac{(\langle B,X\rangle-b)^2}{14(1+b^2)}.
\end{equation}
It is globally $1$-smooth in Frobenius norm, attains its minimum $0$, and satisfies
$f_b(0)=4/7$. Run exact-gradient, exact-polar EF21-SignMuon from
$X_0=M_0=E_0=0$, with either EMA or Nesterov and fixed
$\mu\in[0,1)$.
Let $r_t>0$ satisfy $r_t>\mu r_{t-1}$ for $t\ge2$, and set
\[
 \eta_t=r_t+r_{t+1},\qquad b=-\kappa r_1.
\]
Then, for every $t\ge0$,
\begin{gather*}
 X_t=-hS\sum_{k=1}^t(-1)^{k-1}\eta_k,\\
 \langle B,X_t\rangle-b=(-1)^t\kappa r_{t+1},\\
 \|\nabla f_b(X_t)\|_F\ge1.
\end{gather*}
The same statement holds through any finite horizon when the hypotheses
hold only for the required finite sequence.
\end{theorem}

\begin{paperproof}
The Hessian is the sum of two orthogonal rank-one operators with
eigenvalues $1$ and $1/(1+b^2)$. A minimizer is
$-a_bA+(b/7)B$, and
$f_b(0)=a_b^2/2+b^2/[14(1+b^2)]=4/7$.
In the article's spectral/nuclear smoothness convention, it also satisfies
$\|\nabla f_b(X)-\nabla f_b(Y)\|_*\le8\|X-Y\|_{2\to2}$,
since $\|A\|_*=1$ and $\|B\|_*=7$.
Suppose the claimed trajectory holds through $X_{t-1}$. Its gradient is
\[
 G_t=a_bA+(-1)^{t-1}\frac{\kappa r_t}{7(1+b^2)}B.
\]
The EMA coefficient of $A$ is $a_b(1-\mu^t)>0$.
After removing the sign $(-1)^{t-1}$ and a positive common factor,
its coefficient of $B$ is
\[
 r_t-\mu r_{t-1}+\mu^2r_{t-2}-\cdots+(-\mu)^{t-1}r_1>0.
\]
Pair adjacent terms; each pair is positive by the hypothesis, as is
the unpaired last term. Consecutive EMA coefficients of $B$ therefore
have opposite signs. Nesterov's signal
$(1+\mu)M_t-\mu M_{t-1}$ has the same $B$ sign and a positive $A$
coefficient as well. Lemma~\ref{add:lem:invisible} gives $D_t$ and $E_t$.
The update $X_t=X_{t-1}-\eta_tE_t$ and
$\eta_t=r_t+r_{t+1}$ complete the induction.
Finally, $\langle A,X_t\rangle=0$ implies
$\langle A,\nabla f_b(X_t)\rangle=a_b\ge1$, proving the gradient floor.
\end{paperproof}

\paragraph{Decreasing steps.}
With $\mu=0$, every strictly decreasing positive schedule
$\eta_t\to0$ is covered: take
$r_t=\sum_{k=t}^{\infty}(-1)^{k-t}\eta_k$.
The alternating-series bounds give $0<r_t<\eta_t$ and
$r_t+r_{t+1}=\eta_t$. Thus one fixed quadratic defeats the entire run,
including $\eta_t=a/t$ and $a/\sqrt t$. The iterates converge to a point
whose gradient is $a_bA$, of norm at least $1$.
For any fixed $\mu<1$, choose $a>0$, $0<p\le1$, and $\tau\ge0$
large enough that
\[
 r_t=\frac{a}{(t+\tau)^p},\qquad
 \left(\frac{1+\tau}{2+\tau}\right)^p>\mu.
\]
Both momentum variants are then covered. These steps satisfy
$\eta_t\asymp t^{-p}$ and
$\sum_t\eta_t=\infty$; choosing $p>1/2$ also gives
$\sum_t\eta_t^2<\infty$.

\paragraph{Parameters chosen for a finite budget.}
For every horizon $T$, any constant $\eta_T>0$ and any fixed
$\mu_T\in[0,1)$, choose $r_t=\eta_T/2$ and
$b_T=-\kappa\eta_T/2$. Then
\begin{gather*}
 \min_{0\le t\le T}\|\nabla f_{b_T}(X_t)\|_F\ge1,\\
 L_F=1,\qquad f_{b_T}(0)-\inf f_{b_T}=4/7.
\end{gather*}
Thus no uniform vanishing stationarity rate over this smooth,
lower-bounded class follows by tuning a constant learning rate and momentum
to $T$. The dimension and compressor are also fixed, and the noise is zero.
With $\mu=0$, any finite nonincreasing positive schedule is covered:
set $r_{T+1}=\eta_T/2$ and recursively $r_t=\eta_t-r_{t+1}>0$.
These are uniform worst-case statements, with the instance allowed to
depend on the chosen schedule; arbitrary varying momentum is not asserted.

\paragraph{Why contraction does not suffice.}
Here $\kappa\simeq5.039$. The achieved scaled-sign contraction is
$\kappa/8\simeq0.630$ initially and
$4\kappa/(8+3\kappa)\simeq0.872$ thereafter.
Nevertheless, every estimator is orthogonal to the fixed gradient
component $a_bA$. All finite-time momentum signals have full rank and
all compressed residuals have strict signs.
As the steps decrease, the gradients approach rank one
while their polar factors keep alternating. Smoothness of the objective
therefore does not make the polar target slowly varying.

\paragraph{Exact algebraic certificate.}
The following data make Lemma~\ref{add:lem:invisible} explicit.
Let $K$ be the skew-symmetric integer matrix
\[
\setlength{\arraycolsep}{2pt}
K=\begin{pmatrix}
0&-6&-10&15&4&-3&-9&3\\
6&0&-3&-7&0&-12&-4&0\\
10&3&0&20&0&8&-1&4\\
-15&7&-20&0&-21&-7&0&-1\\
-4&0&0&21&0&-6&-15&7\\
3&12&-8&7&6&0&-17&0\\
9&4&1&0&15&17&0&12\\
-3&0&-4&1&-7&0&-12&0
\end{pmatrix},
\]
and put $D=(20I-K)(20I+K)^{-1}$, $S=\operatorname{sign}(D)$,
$w=(-1,-8,-9,0,0,-12,7,-8)^\top$.
Let $z=w+\xi e_6$, where $\xi\in(251/250,201/200)$ is the unique root of
\[
3903331846\xi^2+1548933728\xi-5493727089=0.
\]
Define
\[
 A=\frac{Dzz^\top}{z^\top z},\qquad B=D-A.
\]
Skew-symmetry gives $D^\top D=I$; the polynomial is a nonzero
multiple of $z^\top D^\top Sz$, so $\langle A,S\rangle=0$.
The rank-one orthogonal projector $zz^\top/(z^\top z)$ gives the stated
norms and polar identity. Rational interval arithmetic on the displayed
interval verifies, entrywise,
\[
 S_{ij}D_{ij}>1/50,\qquad S_{ij}(D-2A)_{ij}>1/50.
\]
For reproducibility, the accompanying script
\path{research/ef21_signmuon_vanishing_steps.py} verifies these
inequalities with exact fractions: after multiplying the second by
$z^\top z$, each is a quadratic inequality checked at the interval
endpoints and, when applicable, its vertex. It also checks root existence,
uniqueness, and orthogonality exactly; no numerical search is part of the
certificate.
